% ============================================================
\chapter{Directions de Recherche}
\label{ch:recherche}
% ============================================================

\begin{intuition}
L'apprentissage g\'eom\'etrique est un domaine en pleine expansion.
Ce chapitre pr\'esente les directions de recherche les plus actives
et les probl\`emes ouverts qui d\'efiniront les prochaines avanc\'ees
du domaine.
\end{intuition}

% ============================================================
\section{Mod\`eles de fondation g\'eom\'etriques}
\label{sec:foundation_models}
% ============================================================

\begin{definition}[Mod\`ele de fondation g\'eom\'etrique]
\index{mod\`ele de fondation}
Un \textbf{mod\`ele de fondation g\'eom\'etrique} est un mod\`ele
pr\'e-entra\^\i n\'e \`a grande \'echelle sur des donn\'ees
structur\'ees (graphes, mol\'ecules, nuages de points) et
transf\'erable \`a de multiples t\^aches en aval par
\emph{fine-tuning} ou \emph{prompting}.
\end{definition}

\begin{exemple}[Exemples de mod\`eles de fondation g\'eom\'etriques]
\begin{itemize}
  \item \textbf{GEM} (Fang et al.\ 2022)~: mod\`ele de fondation
        pour les mol\'ecules, pr\'e-entra\^\i n\'e sur 20M de
        structures.
  \item \textbf{Uni-Mol} (Zhou et al.\ 2023)~: repr\'esentation
        unifi\'ee 2D/3D pour les mol\'ecules.
  \item \textbf{Point-MAE} (Pang et al.\ 2022)~: auto-encodeur
        masqu\'e pour les nuages de points 3D.
\end{itemize}
\end{exemple}

\begin{remarque}
Contrairement aux mod\`eles de fondation textuels (GPT, Claude),
les mod\`eles g\'eom\'etriques doivent respecter les
\textbf{sym\'etries} des donn\'ees (invariance par permutation,
\'equivariance par rotation). Cela impose des contraintes
architecturales fortes.
\end{remarque}

\begin{proposition}[D\'efis du scaling]
Le passage \`a l'\'echelle des GNN pose des probl\`emes sp\'ecifiques~:
\begin{enumerate}
  \item \textbf{Explosion du voisinage}~: $k$ couches impliquent
        $\mathcal{O}(d^k)$ voisins ($d$ = degr\'e moyen).
  \item \textbf{Over-smoothing}~: les repr\'esentations convergent
        avec la profondeur.
  \item \textbf{Over-squashing}~: l'information longue distance est
        compress\'ee dans des bottlenecks.
  \item \textbf{Graphes de taille variable}~: le batching est
        non trivial.
\end{enumerate}
\end{proposition}

% ============================================================
\section{Transformers \'equivariants}
\label{sec:equivariant_transformers}
% ============================================================

\begin{definition}[Transformer g\'eom\'etrique]
\index{Transformer g\'eom\'etrique}
Un \textbf{Transformer g\'eom\'etrique} est un Transformer dont
l'attention int\`egre des informations g\'eom\'etriques tout en
respectant les sym\'etries~:
\[
  \alpha_{ij} = \frac{\exp\!\left(\frac{q_i^\top k_j}{\sqrt{d}}
  + b(x_i, x_j)\right)}{\sum_{l} \exp\!\left(
  \frac{q_i^\top k_l}{\sqrt{d}} + b(x_i, x_l)\right)}
\]
o\`u $b(x_i, x_j)$ est un biais d'attention d\'ependant de la
g\'eom\'etrie (distance, angle, orientation relative).
\end{definition}

\begin{exemple}[Exemples d'architectures]
\begin{itemize}
  \item \textbf{SE(3)-Transformer} (Fuchs et al.\ 2020)~: attention
        $\mathrm{SE}(3)$-\'equivariante avec features tensorielles.
  \item \textbf{Equiformer} (Liao \& Smidt, 2023)~: Transformer
        \'equivariant avec produits tensoriels dans l'attention.
  \item \textbf{GPS} (Ramp\'a\v{s}ek et al.\ 2022)~: combine
        message passing local et attention globale.
\end{itemize}
\end{exemple}

\begin{theoreme}[Expressivit\'e des Graph Transformers]
\index{Graph Transformer}
Un Graph Transformer avec encodage positionnel spectral et attention
globale est \textbf{strictement plus expressif} que le $1$-WL test.
Avec des encodages positionnels $\mathrm{SE}(3)$-\'equivariants,
il peut approcher toute fonction \'equivariante continue.
\end{theoreme}

\begin{formules}[title=\textbf{Attention \'equivariante}]
\[
  h_i' = \sum_{j} \alpha_{ij} \cdot V(h_j, x_j - x_i)
\]
o\`u $V : \R^d \times \R^3 \to \R^{d'}$ est une fonction de valeur
qui d\'epend de la position relative, et $\alpha_{ij}$ est invariant
par $\mathrm{SE}(3)$.
\end{formules}

% ============================================================
\section{Mod\`eles de diffusion sur vari\'et\'es}
\label{sec:diffusion_manifolds}
% ============================================================

\begin{definition}[Diffusion riemannienne]
\index{diffusion riemannienne}
Un \textbf{mod\`ele de diffusion riemannien} (De Bortoli et al.\
2022) g\'en\'eralise les mod\`eles de diffusion (DDPM) aux
vari\'et\'es riemanniennes $(\mathcal{M}, g)$~:
\begin{enumerate}
  \item \textbf{Processus forward}~: mouvement brownien riemannien
        $dX_t = \sqrt{2\beta_t}\, dB_t^{\mathcal{M}}$.
  \item \textbf{Score matching}~: apprentissage du score
        $\nabla_{\mathcal{M}} \log p_t(x)$ sur la vari\'et\'e.
  \item \textbf{Processus reverse}~: int\'egration de l'\'equation
        reverse sur $\mathcal{M}$ avec le score appris.
\end{enumerate}
\end{definition}

\begin{exemple}[Applications de la diffusion g\'eom\'etrique]
\begin{itemize}
  \item \textbf{G\'en\'eration de mol\'ecules 3D}~: diffusion sur
        $\R^{3n}$ avec \'equivariance $\mathrm{SE}(3)$
        (EDM, Hoogeboom et al.\ 2022).
  \item \textbf{G\'en\'eration de conformations}~: diffusion
        sur les angles de torsion (vari\'et\'e torique).
  \item \textbf{G\'en\'eration de prot\'eines}~: diffusion sur
        $\mathrm{SO}(3)^n$ pour les frames des r\'esidus
        (RFDiffusion, Watson et al.\ 2023).
\end{itemize}
\end{exemple}

\begin{attention}[title=\textbf{Difficult\'es techniques}]
La diffusion sur vari\'et\'es n\'ecessite~:
\begin{itemize}
  \item La carte exponentielle $\exp_x : T_x\mathcal{M} \to
        \mathcal{M}$ pour les pas d'int\'egration.
  \item Le logarithme riemannien $\log_x : \mathcal{M} \to
        T_x\mathcal{M}$ pour le score.
  \item Le transport parall\`ele pour d\'eplacer les vecteurs
        entre espaces tangents.
\end{itemize}
\end{attention}

% ============================================================
\section{Apprentissage auto-supervis\'e g\'eom\'etrique}
\label{sec:ssl_geometrique}
% ============================================================

\begin{definition}[Apprentissage auto-supervis\'e sur graphes]
\index{apprentissage auto-supervis\'e}
L'\textbf{apprentissage auto-supervis\'e} (SSL) sur graphes
apprend des repr\'esentations sans labels via des t\^aches
auxiliaires~:
\begin{enumerate}
  \item \textbf{Contrastif}~: maximiser la similarit\'e entre
        deux vues augment\'ees du m\^eme graphe (GraphCL, GCA).
  \item \textbf{Pr\'edictif}~: pr\'edire des propri\'et\'es
        masqu\'ees (features de n\oe uds, ar\^etes, sous-graphes).
  \item \textbf{G\'en\'eratif}~: reconstruire le graphe \`a
        partir d'une repr\'esentation compress\'ee (GraphMAE).
\end{enumerate}
\end{definition}

\begin{proposition}[Augmentations g\'eom\'etriques]
Les augmentations pour les graphes incluent~:
\begin{itemize}
  \item Suppression al\'eatoire de n\oe uds ou d'ar\^etes.
  \item Perturbation des features.
  \item Sous-graphe al\'eatoire.
  \item Diffusion de chaleur (lissage des features).
\end{itemize}
Pour les donn\'ees 3D, on ajoute~: rotation, translation,
bruit gaussien, \'echantillonnage partiel.
\end{proposition}

\begin{theoreme}[Garanties du SSL contrastif]
\index{apprentissage contrastif}
Sous des hypoth\`eses de r\'egularit\'e sur la distribution des
donn\'ees et des augmentations, les repr\'esentations contrastives
sont \textbf{lin\'eairement s\'eparables} pour les t\^aches en
aval, avec une borne d'erreur contr\^ol\'ee par l'alignement et
l'uniformit\'e des repr\'esentations.
\end{theoreme}

% ============================================================
\section{Probl\`emes ouverts}
\label{sec:problemes_ouverts}
% ============================================================

\begin{remarque}[Probl\`emes ouverts majeurs]
\index{probl\`emes ouverts}
\begin{enumerate}
  \item \textbf{Expressivit\'e optimale}~: quelle est l'architecture
        GNN minimale pour \'egaler $k$-WL pour un $k$ donn\'e~?
  \item \textbf{Over-squashing}~: comment permettre l'information
        longue distance sans attention globale co\^uteuse
        ($\mathcal{O}(n^2)$)~?
  \item \textbf{G\'en\'eralisation hors distribution}~: les GNN
        g\'en\'eralisent-ils \`a des graphes de structures tr\`es
        diff\'erentes de l'entra\^\i nement~?
  \item \textbf{\'Equivariance approch\'ee}~: comment g\'erer les
        sym\'etries approximatives (quasi-sym\'etries) des donn\'ees
        r\'eelles~?
  \item \textbf{Interpr\'etabilit\'e}~: comment expliquer les
        d\'ecisions d'un GNN en termes de sous-structures
        significatives~?
  \item \textbf{Passage \`a l'\'echelle}~: comment entra\^\i ner
        des GNN sur des graphes \`a des milliards de n\oe uds~?
\end{enumerate}
\end{remarque}

\begin{definition}[Over-squashing]
\index{over-squashing}
L'\textbf{over-squashing} d\'esigne le ph\'enom\`ene o\`u
l'information provenant de n\oe uds distants est compress\'ee
de mani\`ere exponentielle en traversant les couches de message
passing. Formellement, le jacobien~:
\[
  \left\|\frac{\partial h_v^{(L)}}{\partial h_u^{(0)}}\right\|
  \leq C \cdot \lambda_{\max}^L \cdot
  \prod_{\ell=1}^L \mathrm{Lip}(\sigma_\ell)
\]
d\'ecro\^\i t exponentiellement avec la distance dans le graphe
si $\lambda_{\max} < 1$.
\end{definition}

% ============================================================
\section{Exercices}
% ============================================================

\begin{exercice}
Expliquer pourquoi les encodages positionnels laplaciens ne sont
pas uniques (signe, ordre). Proposer une solution pour les rendre
invariants.
\end{exercice}

\begin{exercice}
Pour un graphe mol\'eculaire, proposer trois augmentations
contrastives qui pr\'eservent l'identit\'e chimique et trois qui
ne la pr\'eservent pas. Justifier.
\end{exercice}

\begin{exercice}[Over-squashing]
Calculer la borne sur $\left\|\frac{\partial h_v^{(L)}}
{\partial h_u^{(0)}}\right\|$ pour un graphe en chemin de
longueur $L$ avec un GCN \`a poids unitaires et activation ReLU.
\end{exercice}

\begin{exercice}
Montrer que le mouvement brownien sur la sph\`ere $S^2$ converge
vers la distribution uniforme. Quel est le taux de convergence
en fonction de la premi\`ere valeur propre non nulle du laplacien~?
\end{exercice}

\begin{exercice}[Projet]
Impl\'ementer un mod\`ele de diffusion \'equivariant simple pour
g\'en\'erer des nuages de points 3D (par exemple, des formes de
la base ShapeNet). \'Evaluer la qualit\'e par la Chamfer distance
et les nombres de Betti.
\end{exercice}
